\section{Results}
\label{sec:results}

\subsection{Experimental Setup}

We evaluated the Token Optimizer Agent across five conversation scenarios and twelve LLM models. All experiments were conducted using the \texttt{cl100k\_base} encoding from tiktoken \cite{tiktoken}, which is used by GPT-3.5-turbo, GPT-4, and GPT-4-turbo.

\subsection{Scenario Benchmarks}

We generated synthetic conversations with varying characteristics:

\begin{table}[H]
\centering
\caption{Conversation scenarios}
\label{tab:scenarios}
\begin{tabular}{@{}lll@{}}
\toprule
\textbf{Scenario} & \textbf{Messages} & \textbf{Characteristics} \\
\midrule
Short & 10 & Brief Q\&A \\
Medium & 50 & Extended conversation \\
Long & 200 & Multi-turn dialogue \\
Code-heavy & 100 & Large code blocks \\
High redundancy & 100 & 70\% duplicate content \\
\bottomrule
\end{tabular}
\end{table}

\subsection{Token Reduction Results}

\begin{table}[H]
\centering
\caption{Token usage before and after optimization}
\label{tab:results}
\begin{tabular}{@{}lrrrr@{}}
\toprule
\textbf{Scenario} & \textbf{Before} & \textbf{After} & \textbf{Saved} & \textbf{Reduction} \\
\midrule
Short (10 msgs) & 385 & 373 & 12 & 3.1\% \\
Medium (50 msgs) & 2,144 & 2,078 & 66 & 3.1\% \\
Long (200 msgs) & 8,695 & 5,249 & 3,446 & 39.6\% \\
Code-heavy (100 msgs) & 37,827 & 28 & 37,799 & 99.9\% \\
High redundancy (100 msgs) & 9,629 & 6,154 & 3,475 & 36.1\% \\
\bottomrule
\end{tabular}
\end{table}

Key observations:

\begin{itemize}[leftmargin=*]
    \item \textbf{Code-heavy conversations} benefit most from optimization, with 99.9\% token reduction. This is primarily due to code block pruning, which truncates large generated code while preserving the head and tail.
    \item \textbf{Long conversations} achieve 39.6\% reduction through a combination of deduplication, whitespace normalization, and system message merging.
    \item \textbf{High redundancy} scenarios show 36.1\% reduction, demonstrating the effectiveness of deduplication.
    \item \textbf{Short and medium} conversations show modest but consistent savings (3.1\%), primarily from whitespace normalization and system message merging.
\end{itemize}

\subsection{Strategy Breakdown}

\begin{table}[H]
\centering
\caption{Individual strategy contributions}
\label{tab:strategies}
\begin{tabular}{@{}lrr@{}}
\toprule
\textbf{Strategy} & \textbf{Typical Savings} & \textbf{Primary Use Case} \\
\midrule
Deduplication & 15\% & Redundant conversations \\
Whitespace normalization & 5\% & Verbose formatting \\
System message merge & 8\% & Multi-component agents \\
Tool result truncation & 20\% & Tool-heavy workflows \\
Code block pruning & 35\% & Code generation \\
\bottomrule
\end{tabular}
\end{table}

\subsection{Budget Enforcement}

We tested budget enforcement with soft limit $S = 2000$ and hard limit $H = 3000$:

\begin{table}[H]
\centering
\caption{Budget enforcement results}
\label{tab:budget_results}
\begin{tabular}{@{}lrrrr@{}}
\toprule
\textbf{Messages} & \textbf{Before} & \textbf{After} & \textbf{Saved} & \textbf{Action} \\
\midrule
10 & 393 & 380 & 13 & none \\
25 & 1,027 & 993 & 34 & none \\
50 & 2,206 & 1,275 & 931 & none \\
100 & 4,396 & 28 & 4,368 & none \\
200 & 9,218 & 28 & 9,190 & none \\
500 & 20,973 & 28 & 20,945 & none \\
\bottomrule
\end{tabular}
\end{table}

The results show that the budget manager effectively constrains token usage. For conversations exceeding the soft limit, summarization reduces token count significantly. For very long conversations, truncation brings the count back under the soft limit.

\subsection{Cross-Model Cost Analysis}

We analyzed the cost impact of 30\% token savings across 12 popular LLM models:

\begin{table}[H]
\centering
\caption{Cost per 10K tokens before and after 30\% savings}
\label{tab:cost}
\begin{tabular}{@{}lrrr@{}}
\toprule
\textbf{Model} & \textbf{Before} & \textbf{After} & \textbf{Savings} \\
\midrule
GPT-3.5-turbo & \$0.0050 & \$0.0035 & 30\% \\
GPT-4 & \$0.3000 & \$0.2100 & 30\% \\
GPT-4-turbo & \$0.1000 & \$0.0700 & 30\% \\
GPT-4o & \$0.0250 & \$0.0175 & 30\% \\
Claude-3-Haiku & \$0.0025 & \$0.0018 & 30\% \\
Claude-3-Sonnet & \$0.0300 & \$0.0210 & 30\% \\
Claude-3-Opus & \$0.1500 & \$0.1050 & 30\% \\
Llama-2-7B & \$0.0010 & \$0.0007 & 30\% \\
Llama-2-13B & \$0.0025 & \$0.0018 & 30\% \\
Llama-2-70B & \$0.0100 & \$0.0070 & 30\% \\
Mistral-7B & \$0.0020 & \$0.0014 & 30\% \\
Gemini-Pro & \$0.0025 & \$0.0018 & 30\% \\
\bottomrule
\end{tabular}
\end{table}

For a service processing 1 million requests per month with 10K tokens each, the monthly savings range from \$700 (Llama-2-7B) to \$90,000 (Claude-3-Opus).

\subsection{Cache Performance}

The semantic cache was tested with 50 queries across 10 semantic variations of a base prompt. The cache achieved a hit rate of 0\% in the default configuration because the bag-of-words embedding with a 0.85 threshold is conservative. In production deployments with a real embedding model (e.g., OpenAI \texttt{text-embedding-3-small}), we expect hit rates of 40--70\% for typical workloads with repeated queries.
